% =========================================================
% PACOTES FUNDAMENTAIS E IDIOMA
% =========================================================
\usepackage[T1]{fontenc}        % Seleção de códigos de fonte.
\usepackage[utf8]{inputenc}     % Codificação do documento (conversão automática de acentos)
\usepackage[brazil]{babel}      % Idioma do documento (traduções de "Chapter" para "Capítulo")
\usepackage{lmodern}            % Usa a fonte Latin Modern (alta qualidade em PDF)
\usepackage{textcomp}           % Símbolos complementares
\usepackage{bold-extra}         % Permite negrito e itálico/small-caps simultâneos simultaneamente

% =========================================================
% ESTRUTURA, MARGENS E LAYOUT (Padrão ABNT)
% =========================================================
% Configuração cravada de margens ABNT (3cm esq/sup, 2cm dir/inf)
\usepackage[left=3cm,top=3cm,right=2cm,bottom=2cm]{geometry}
\usepackage{setspace}           % Permite o comando \onehalfspacing (espaçamento 1.5 padrão ABNT)
\usepackage{indentfirst}        % Indenta obrigatoriamente o primeiro parágrafo de cada seção
\usepackage{lastpage}           % Conta o total de páginas (usado em fichas catalográficas e rodapés)
\usepackage{multicol} % Permite usar o ambiente \begin{multicols}


% =========================================================
% TABELAS E PÁGINAS COMPLEXAS
% =========================================================
\usepackage{longtable}          % Tabelas que quebram em várias páginas (essencial para roteiros)
\usepackage{tabularx}           % Tabelas com largura de coluna auto-ajustável (\begin{tabularx}{\textwidth})
\usepackage{makecell}           % Permite quebras de linha manuais dentro de células de tabela
\usepackage{pdflscape}          % Permite girar páginas específicas para modo paisagem
\usepackage{eso-pic}            % Inserção de imagens de fundo (marcas d'água, capas customizadas)

% =========================================================
% GRÁFICOS, CORES E FLUTUANTES
% =========================================================
\usepackage{graphicx}           % Inclusão de imagens e gráficos
\usepackage{subcaption}         % Permite subfiguras (1a, 1b, etc)
\usepackage{xcolor}             % Suporte a cores mais avançado que o antigo pacote 'color'
\usepackage{float}              % Controle rígido de flutuantes (permite o uso do [H] para fixar tabelas/figuras)
\usepackage{placeins}           % Fornece o \FloatBarrier para impedir que imagens invadam outras seções

% =========================================================
% UNIDADES E MATEMÁTICA
% =========================================================
\usepackage{units}              % Formatação padronizada de unidades de medida
% =========================================================
% PACOTES DE MATEMÁTICA
% =========================================================
\usepackage{amsmath}    % Resolve o seu erro atual (ambientes matemáticos)
\usepackage{amsfonts}   % Fontes matemáticas avançadas
\usepackage{amssymb}    % Símbolos matemáticos adicionais

% =========================================================
% MICROTIPOGRAFIA E AJUSTES FINOS
% =========================================================
\usepackage{microtype}          % Resolve 90% dos problemas de justificação (overfull/underfull hbox)
\setlength{\emergencystretch}{2em} % Dá uma folga extra para o LaTeX não jogar palavras para fora da margem
\setlength{\parindent}{1.25cm}  % Define o recuo exato do parágrafo no padrão ABNT (1.25cm a 1.5cm)

% =========================================================
% CITAÇÕES ABNT E REFERÊNCIAS
% =========================================================
% Backref DEVE vir antes do hyperref (mostra em quais páginas a referência foi citada)
\usepackage[brazilian,hyperpageref]{backref} 
% Pacote oficial de citações ABNT em formato alfabético (AUTOR, Ano)
\usepackage[alf]{abntex2cite}   

% =========================================================
% HYPERLINKS (SEMPRE O ÚLTIMO PACOTE A SER CARREGADO)
% =========================================================
\usepackage{hyperref}           % Cria links clicáveis no PDF
\hypersetup{
    colorlinks=true,            % false = links em caixas vermelhas, true = links coloridos
    linkcolor=black,            % Cor dos links internos (Sumário, figuras, tabelas) - Preto para impressão/acadêmico
    citecolor=black,            % Cor dos links de citações [1] ou (AUTOR, 2026)
    urlcolor=blue,              % Cor dos links de internet (\url)
    bookmarksdepth=4            % Nível de profundidade dos marcadores laterais no leitor de PDF
}